\documentclass[12pt]{article}
\usepackage[utf8]{inputenc}
\usepackage[T1]{fontenc}
\usepackage{hyperref}
\title{Blissymbols Language Resources for K-12 Curriculum Development}
\author{Compiled by Claude Sonnet 5.5 (claude-sonnet-5-5) for Andr\'as Kornai}
\date{}
\begin{document}
\maketitle

\section{Introduction}
Blissymbolics (Blissymbols, ISO 639-3 zbl) is a constructed, graphic, semantically based language devised by Charles K. Bliss. It was first used functionally at the Ontario Crippled Children's Centre in Toronto in 1971. It is used mainly for augmentative and alternative communication and for literacy support. It has no native-speaking community in the usual sense, so ``native'' materials here means materials produced by the Blissymbolics community itself. Most sources found are older (1980s--1990s), and I could not confirm whether their print items are still available.

\section{Primary Educational Resources}
\begin{itemize}
\item Blissymbolics Communication International (BCI) training resources, for example the web training page: \url{https://www.blissymbolics.org/WebTraining/NewResourceHTMLCSS/resource9.php}
\item Print titles listed in an educational technology bibliography: \emph{Blissymbol Reference Guide} (BCI, 1992), \emph{Teaching and Using Blissymbolics} (McDonald, 1980) and \emph{Communicating with Blissymbolics} (McNaughton, ed., 1985): \url{https://edutechwiki.unige.ch/en/Blissymbols}
\end{itemize}

\section{Linguistic and Reference Materials}
\begin{itemize}
\item University of Toronto OISE library guide to Blissymbolics: \url{https://guides.library.utoronto.ca/oise-ohec/blissymbolics}
\item Evertype page on Blissymbols standards (encoding and standards context): \url{https://evertype.com/standards/by/bliss.html}
\item Blissymbolics Communication Institute Canada archival fonds, University of Toronto, covering the development, dissemination and use of Blissymbols: \url{https://discoverarchives.library.utoronto.ca/index.php/blissymbolics-communication-institute-canada-fonds}
\item Archive description of the Blissymbolics Communication Institute Canada: \url{https://discoverarchives.library.utoronto.ca/index.php/blissymbolics-communication-institute-canada}
\end{itemize}

\section{Religious and Community Resources}
No religious materials turned up in the searches.
\begin{itemize}
\item Blissymbolics Communication International, a non-profit charity that holds the exclusive license to publish Blissymbols for people with communication, language and learning difficulties: \url{https://www.canadahelps.org/en/charities/blissymbolics-communication-international}
\end{itemize}

\section{Digital and Audio Resources}
No audio resources were found. The BCI web training pages listed above are the main online material found.

\section{Curriculum Development Context}
Blissymbols are taught mainly to non-speaking learners as a route to communication and literacy, not as a mother tongue in schools. A K-12 plan would therefore be an AAC and literacy programme, built on the BCI vocabulary and training materials above. The older print texts would need checking for availability. The Blissymbolics Learning Centre in Muskoka was named in search results as offering distance education, but I did not verify that it is still running.

\end{document}